\section{Zadanie 8 (Oskar Prusik)}
Niech \( a, b \) będą liczbami rzeczywistymi. Jakie warunki muszą spełniać \( a \) i \( b \), żeby istniała zmienna losowa, która ma wartość oczekiwaną równą \( a \) i wariancję równą \( b \)?  
W przypadku, gdy taka zmienna losowa istnieje, podaj ją.

\subsection*{Rozwiązanie}
Oczywistym jest, że wariancja jest nieujemna, co nakłada nam juz warunek: \[b \geq 0\]

Teraz znajdujemy przykład zmiennej losowej dla tylko tego jednego warunku. Niech zmienna losowa \(X\) ma rozkład:
\[
P(X = a + \sqrt{b}) = \frac{1}{2}, \quad P(X = a - \sqrt{b}) = \frac{1}{2}
\]
gdzie \(a \in \mathbb{R}\), \(b \geq 0\).


Korzystając z definicji wartości oczekiwanej dla zmiennej dyskretnej:
\[
E[X] = \sum x_i \cdot p_i = \frac{1}{2} \cdot (a + \sqrt{b}) + \frac{1}{2} \cdot (a - \sqrt{b})
\]
\[
= \frac{1}{2}a + \frac{1}{2}\sqrt{b} + \frac{1}{2}a - \frac{1}{2}\sqrt{b} = a
\]


\[
E[X^2] = \frac{1}{2} \cdot (a + \sqrt{b})^2 + \frac{1}{2} \cdot (a - \sqrt{b})^2
\]
\[
= \frac{1}{2} \cdot (a^2 + 2a\sqrt{b} + b) + \frac{1}{2} \cdot (a^2 - 2a\sqrt{b} + b)
\]
\[
= \frac{1}{2}a^2 + a\sqrt{b} + \frac{1}{2}b + \frac{1}{2}a^2 - a\sqrt{b} + \frac{1}{2}b
\]
\[
= a^2 + b
\]


Korzystając ze wzoru \(\mathrm{Var}(X) = E[X^2] - (E[X])^2\):
\[
\mathrm{Var}(X) = (a^2 + b) - a^2 = b
\]


\subsection*{Odpowiedź}
\[
\boxed{\text{Dla dowolnych \(a \in \mathbb{R}\) i \(b \geq 0\) istnieje zmienna losowa o podanej wartości oczekiwanej i wariancji.}}
\]
